
\subsection{Thủ thuật vùng chi dưới}

\begin{longtblr}[
    caption = {Thủ thuật xoa bóp vùng chi dưới},
]{
    width = \textwidth,
    colspec = {| X[1,c] |X[1, l]|X[2, c]|X[1, c]|X[7, l]|},
    % hlines,
    rowhead = 1,
}
    \hline
    ID & Động tác & Vùng tác động & Số lần & Nội dung 
    \\ \hline
    1 & Xoa & & & {%
        Từ gốc chi đi xuống
    }
    \\ \hline
    2 & Xát & & & {%
        Từ gốc chi đi xuống
    }
    \\ \hline
    3 & Miết & & & {%
        Từ trên xuống
    }
    \\ \hline
    4 & Phân hợp & & & {%
        Từ trên xuống
    }
    \\ \hline
    5 & Vỗ & & & {%
    }
    \\ \hline
    6 & Day & & & {%
    }
    \\ \hline
    7 & Đấm & & & {%
    }
    \\ \hline
    8 & Chặt & & & {%
    }
    \\ \hline
    9 & Lăn & & & {%
    }
    \\ \hline
    10 & Bóp & & & {%
    }
    \\ \hline
    11 & Vờn & & & {%
    }
    \\ \hline
    12 & Bấm & Huyệt & 1 & {%
        Dũng Tuyền,
        Côn Lôn,
        Thái Khê,
        Huyền Chung,
        Tam Âm Giao,
        Thừa Sơn,
        Uỷ Trung,
        Ân Môn,
        Thừa Phủ,
        Trật Biên (gấp gót, vuông góc, mũi gót ở đâu thì huyệt ở đó),
        Phong Thị(vai để tự do, đầu ngón giữa ở đâu huyệt ở đó),
        Dương Lăng Tuyền,
        Túc Tam Lý,
        Bát Phong,
        Thái Xung,
        Nội Đình,
        Giải Khê,
        Độc Tỵ,
        Tất Nhãn,
        Hạc Đỉnh,
        Huyết Hải,
        Lương Khâu    
    }
    \\ \cline{5-5} & & & & {%
        Dũng Tuyền

        \dungtuyen
    }
    \\ & & & & {%
        Côn Lôn

        \conlon
    }
    \\ & & & & {%
        Thái Khê

        \thaikhe
    }
    \\ & & & & {%
        Huyền Chung

        \huyenchung
    }
    \\ & & & & {%
        Tam Âm Giao

        \tamamgiao
    }
    \\ & & & & {%
        Thừa Sơn

        \thuason
    }
    \\ & & & & {%
        Uỷ Trung

        \uytrung
    }
    \\ & & & & {%
        Ân Môn

        \anmon
    }
    \\ & & & & {%
        Thừa Phủ

        \thuaphu
    }
    \\ & & & & {%
        Trật Biên

        \tratbien
    }
    \\ & & & & {%
        Phong Thị

        \phongthi
    }
    \\ & & & & {%
        Dương Lăng Tuyền

        \duonglangtuyen
    }
    \\ & & & & {%
        Túc Tam Lý

        \tuctamly
    }
    \\ & & & & {%
        Bát Phong

        \batphong
    }
    \\ & & & & {%
        Thái Xung

        \thaixung
    }
    \\ & & & & {%
        Nội Đình

        \noidinh
    }
    \\ & & & & {%
        Giải Khê

        \giaikhe
    }
    \\ & & & & {%
        Độc Tỵ

        \docty
    }
    \\ & & & & {%
        Tất Nhãn

        \tatnhan
    }
    \\ & & & & {%
        Hạc Đỉnh

        \hacdinh
    }
    \\ & & & & {%
        Huyết Hải

        \huyethai
    }
    \\ & & & & {%
        Lương Khâu

        \luongkhau
    }
    \\ \hline
    12 & Vận động & Chi dưới & & {%
        % khớp háng (ngửa, mở khớp háng, giữ khớp gối và xương chậu, nhún nhả liên tục, nhún căng nhả) \newline
        % 15. Đẩy đầu gối và khớp háng (sau đó xoay và mở khớp háng, đẩy là tịnh tiến, xoay là vòng) \newline
        % 16. mở 120, thay vì đẩy thì ép. \newline

        Vận động khớp háng mở khớp háng.

        \hyperlink{target:day_dau_goi}{\textbf{Vận động khớp gối tư thế nằm ngửa}}

        \hyperlink{target:xoay_co_chan}{\textbf{Xoay cổ chân tư thế nằm ngửa}}

        \hyperlink{target:gap_duoi_co_chan}{\textbf{Gấp duỗi cổ chân tư thế nằm ngửa}}

        \hyperlink{target:nghieng_co_chan}{\textbf{Nghiêng cổ chân tư thế nằm ngửa}}
        
        \hyperlink{target:ve_ngon_chan}{\textbf{Vê ngón chân tư thế nằm ngửa}}

        \hyperlink{target:vuot_ngon_chan}{\textbf{Vuốt ngón chân tư thế nằm ngửa}}

        \hyperlink{target:van_dong_khop_hang_tu_the_nam_ngua}{\textbf{Vận động khớp háng tư thế nằm ngửa}}

        \hyperlink{target:van_dong_khop_hang_tu_the_nam_sap}{\textbf{Vận động khớp háng tư thế nằm sấp}}

        \hyperlink{target:van_dong_khop_goi_tu_the_nam_sap}{\textbf{Vận động khớp gối tư thế nằm sấp}}
    }
    \\ \hline
\end{longtblr}
